% 《分析学》三卷共用的基础排版与语义宏。
% 各卷主文件须先定义 \BookTitle、\BookVolumeName 和 \BookDraftDate。
\newcommand{\BookAuthor}{R. EverStarine}

\usepackage[a4paper,inner=28mm,outer=24mm,top=27mm,bottom=28mm,
  headheight=16pt,headsep=8mm,footskip=12mm]{geometry}

\usepackage{amsmath}
\usepackage{mathtools}
\usepackage{amsthm}
\usepackage{aliascnt}
\usepackage{unicode-math}
\setmainfont{texgyrepagella-regular.otf}[
  BoldFont=texgyrepagella-bold.otf,
  ItalicFont=texgyrepagella-italic.otf,
  BoldItalicFont=texgyrepagella-bolditalic.otf
]
\setsansfont{texgyreheros-regular.otf}[
  BoldFont=texgyreheros-bold.otf,
  ItalicFont=texgyreheros-italic.otf,
  BoldItalicFont=texgyreheros-bolditalic.otf
]
\setmonofont{lmmonolt10-regular.otf}[
  BoldFont=lmmonolt10-bold.otf,
  ItalicFont=lmmonolt10-oblique.otf,
  BoldItalicFont=lmmonolt10-boldoblique.otf
]
% 行内、行间及标题内公式均使用 unicode-math 默认 Latin Modern Math。
% 公式之外的正文西文仍由 setmainfont 使用 Pagella；黑板粗体单独采用 AMS。
\let\mathbb\relax
\DeclareMathAlphabet{\mathbb}{U}{msb}{m}{n}
% unicode-math 的 Latin Modern Math 将 \varnothing 映为瘦长椭圆；空集单独
% 采用传统 AMS MSBM 的圆形斜杠字形，不改变其余数学符号的字体。
\DeclareSymbolFont{AnalysisAMSb}{U}{msb}{m}{n}
\DeclareMathSymbol{\AnalysisRoundEmptySet}{\mathord}{AnalysisAMSb}{"3F}
\AtBeginDocument{%
  \let\varnothing\AnalysisRoundEmptySet
  \let\emptyset\AnalysisRoundEmptySet
}

\usepackage{microtype}
\usepackage{xcolor}
\usepackage{graphicx}
\usepackage{booktabs}
\usepackage{longtable}
\usepackage{array}
\usepackage{tabularx}
\usepackage{algorithm}
\usepackage{algpseudocode}
\usepackage{enumitem}
\usepackage{fancyhdr}
\usepackage[most]{tcolorbox}
\usepackage{tikz-cd}
\usepackage{csquotes}
\usepackage[
  backend=biber,
  style=numeric-comp,
  sorting=nyt,
  giveninits=true,
  maxbibnames=99,
  doi=true,
  url=true,
  isbn=true
]{biblatex}
\addbibresource{../Shared/References.bib}
% 访问日期保留在书目库供内部核验，成品参考文献不显示 visited on。
\renewbibmacro*{url+urldate}{\usebibmacro{url}}
% biblatex 在书目内重绑定格式宏；在其作用域中将强调统一为正体，
% 同时避免书名中的中文及标点仍继承旧的斜体样式。
\AtBeginBibliography{\let\emph\textup\let\textit\textup}
% 中文条目的责任者连接符：不可套用西文的 “and”，按 GB/T 7714—2015 的著录
% 惯例改用逗号分隔。中文条目在 Shared/References.bib 中以 langid = {chinese}
% 标注（含译本，因其著录信息为中文形式），西文条目保持 “and” 不变。
% 按条目切换须在每条书目项处理时判定；\DeclareDelimFormat 的可选参数不是
% 语言选择器（实测无效），故用 \AtEveryBibitem 配合 \iffieldequalstr。
\AtEveryBibitem{%
  \iffieldequalstr{langid}{chinese}%
    {\DeclareDelimFormat{finalnamedelim}{\addcomma\space}%
     \DeclareDelimFormat{multinamedelim}{\addcomma\space}}%
    {\DeclareDelimFormat{finalnamedelim}{\addspace\bibstring{and}\space}%
     \DeclareDelimFormat{multinamedelim}{\addcomma\space}}%
}
\usepackage{etoolbox}
\usepackage{xparse}
\usepackage[noautomatic]{imakeidx}

\definecolor{AnalysisBlue}{HTML}{244A64}
\definecolor{AnalysisBlueLight}{HTML}{EEF4F7}
\definecolor{AnalysisGray}{HTML}{4F5962}
\definecolor{AnalysisGrayLight}{HTML}{F4F5F6}

\ctexset{
  part = {
    % 编题页分两行：第一行“第X编”，与封面“第一卷”同款（大号无衬线灰色）；
    % 第二行编名，与封面“分析学”同款（大号无衬线蓝色粗体），字号保持原样。
    % 第一行与行距由 number／aftername 给出，format 只负责编名的字体，
    % 以免字号被重置或被重复输出。
    format = {\centering\fontsize{30}{38}\selectfont\sffamily\bfseries\color{AnalysisBlue}},
    name = {,},
    number = {\centering\normalfont\normalsize{\fontsize{24}{30}\selectfont\sffamily\color{AnalysisGray}第\chinese{part}编}},
    aftername = {\par\vspace{0.7cm}},
    beforeskip = 0pt,
    afterskip = 2.5\baselineskip,
    pagestyle = plain
  },
  chapter = {
    format = \raggedright\Huge\bfseries\color{AnalysisBlue},
    name = {第,章},
    number = \chinese{chapter},
    aftername = \quad,
    beforeskip = 0pt,
    afterskip = 30pt,
    pagestyle = plain
  },
  section = {
    format = \Large\bfseries\color{AnalysisBlue},
    beforeskip = 2.4ex plus .6ex minus .2ex,
    afterskip = 1.1ex plus .2ex
  },
  subsection = {
    format = \large\bfseries\color{AnalysisGray},
    beforeskip = 2ex plus .4ex minus .2ex,
    afterskip = .8ex plus .2ex
  }
}

\setlength{\parindent}{2em}
\setlength{\parskip}{0pt}
\linespread{1.18}
% 仅用于两个连续的普通正文段落；列表、公式、定理、标题等结构边界不用。
% 段间额外留白统一由这一长度控制，避免改动全局 \parskip 后扰乱各类环境。
\newlength{\AnalysisProseParagraphGap}
\AtBeginDocument{%
  \setlength{\AnalysisProseParagraphGap}{0.40\baselineskip}%
}
\NewDocumentCommand{\ProseParagraphBreak}{}{%
  \par\addvspace{\AnalysisProseParagraphGap}%
}
\setcounter{secnumdepth}{3}
\setcounter{tocdepth}{2}
\allowdisplaybreaks[2]
\emergencystretch=2em
\clubpenalty=10000
\widowpenalty=10000
\displaywidowpenalty=10000

\setlist[itemize]{leftmargin=2.2em,itemsep=.35ex,topsep=.6ex}
\setlist[enumerate]{leftmargin=2.5em,itemsep=.35ex,topsep=.6ex,
  font=\normalfont\upshape}
\setlist[enumerate,1]{label=\arabic*.,ref=\arabic*}
\setlist[enumerate,2]{label=\arabic*),ref=\arabic*}
\setlist[enumerate,3]{label=\alph*),ref=\alph*}
\setlist[description]{leftmargin=2.5em,labelindent=0pt,itemsep=.35ex,topsep=.6ex}

% “下列命题等价”专用列表；普通罗列仍使用 enumerate。
\newlist{equivalences}{enumerate}{2}
\setlist[equivalences,1]{leftmargin=2.8em,itemsep=.35ex,topsep=.6ex,
  label=\Roman*.,ref=\Roman*,font=\normalfont\upshape}
\setlist[equivalences,2]{leftmargin=2.8em,itemsep=.35ex,topsep=.6ex,
  label=\roman*),ref=\roman*,font=\normalfont\upshape}

\pagestyle{fancy}
\fancyhf{}
\fancyhead[LE]{\small\color{AnalysisGray}\thepage\quad\leftmark}
\fancyhead[RO]{\small\color{AnalysisGray}\rightmark\quad\thepage}
\renewcommand{\headrulewidth}{0.35pt}
% 保留数学标题的原始大小写，避免 book 类的 \MakeUppercase
% 与 unicode-math 命令在页眉标记中发生冲突。
% 偶数页页眉显示“章序＋章标题”（如“第一章 集合系、可测空间与可测映射”）；
% 章序与章标题之间以 \quad 分隔，间距与奇数页节眉“节号＋节标题”之间一致；
% 附录章的章序由 \CTEXthechapter 按本卷附录字母（如 A）给出。
\renewcommand{\chaptermark}[1]{\markboth{\CTEXthechapter\quad #1}{}}
\renewcommand{\sectionmark}[1]{\markright{\thesection\quad #1}}
% 章末习题不编号，但进入习题页后应覆盖上一节留下的奇数页页眉；
% 同时写入目录项。PDF 书签由 hyperref 对 \section*{习题} 自动生成，
% 不再另行 \pdfbookmark，避免“习题”下再嵌套一个“习题”。
\newcommand{\BookExercises}{%
  \section*{习题}%
  \markright{习题}%
  \addcontentsline{toc}{section}{习题}%
}
\fancypagestyle{plain}{
  \fancyhf{}
  \fancyfoot[C]{\small\color{AnalysisGray}\thepage}
  \renewcommand{\headrulewidth}{0pt}
}

\usepackage[unicode,colorlinks=true,linkcolor=AnalysisBlue,
  citecolor=AnalysisBlue,urlcolor=AnalysisBlue,pdfborder={0 0 0}]{hyperref}
\usepackage{bookmark}

% unicode-math 会调整原始 Unicode 希腊字母的类别码。若把这些字母直接写进
% \texorpdfstring 的书签分支，辅助文件回读时类别码不同，LaTeX 会在每一轮都
% 误报标签发生变化。书签分支统一写成下列可展开宏，既保留希腊字母显示，
% 又使写入和回读的标签文本稳定。
\newcommand{\AnalysisPdfAlpha}{α}
\newcommand{\AnalysisPdfGamma}{Γ}
\newcommand{\AnalysisPdfLambda}{λ}
\newcommand{\AnalysisPdfMu}{μ}
\newcommand{\AnalysisPdfPi}{π}
\newcommand{\AnalysisPdfSigma}{σ}
\newcommand{\AnalysisPdfTheta}{θ}
\newcommand{\AnalysisPdfXi}{ξ}
\newcommand{\AnalysisPdfZeta}{ζ}

\usepackage[nameinlink,noabbrev]{cleveref}

% CTeX 2.5.10 的标题前缀宏带一个参数，而 TeX Live 2025 中 cleveref
% 仍按普通前缀宏展开它们，会把节号写成含 \endcsname 的损坏元数据。
% 改为恒等包装后，CTeX 传入的标题编号保持不变，\ref 与 \cref 均可正常工作。
\makeatletter
\def\p@part#1{#1}
\def\p@chapter#1{#1}
\def\p@section#1{#1}
\def\p@subsection#1{#1}
\def\p@subsubsection#1{#1}
\def\p@paragraph#1{#1}
\def\p@subparagraph#1{#1}
\makeatother
% 非数学西文字母统一正体；不改变括号类别、标点或数学字体。
% 2026-09-11：非数学模式的西文字母统一使用 Pagella 正体，保留粗细区别。
% 保留字母级边界，不重分类数字、括号、标点及中文。
% 此设置也覆盖列表序号字母和术语原文；数学模式保持原样。
\newfontfamily\AnalysisLetterFont{texgyrepagella-regular.otf}[
  BoldFont=texgyrepagella-bold.otf,
  ItalicFont=texgyrepagella-italic.otf,
  BoldItalicFont=texgyrepagella-bolditalic.otf
]
\ExplSyntaxOn
\bool_new:N \l_analysis_text_letters_bool
\bool_set_true:N \l_analysis_text_letters_bool
\cs_new_protected:Npn \AnalysisPreserveMathText
  { \bool_set_false:N \l_analysis_text_letters_bool }
\cs_new_protected:Npn \analysis_letter_group_begin:
  {
    \group_begin:
    \bool_if:NT \l_analysis_text_letters_bool
      { \AnalysisLetterFont \upshape }
  }
% 继承 Default 的全部 xeCJK 转移，保留中西文间距及标点处理；
% 不能用只有开关分组的转移覆盖 Boundary 或 CJK 的原处理。
\xeCJK_new_class:n { AnalysisLatinLetter }
\seq_map_inline:Nn \g__xeCJK_class_seq
  {
    \str_if_eq:nnF {#1} { AnalysisLatinLetter }
      {
        \xeCJK_copy_inter_class_toks:nnnn
          {#1} {AnalysisLatinLetter} {#1} {Default}
        \xeCJK_copy_inter_class_toks:nnnn
          {AnalysisLatinLetter} {#1} {Default} {#1}
        \xeCJK_app_inter_class_toks:nnn {#1} {AnalysisLatinLetter}
          { \analysis_letter_group_begin: }
        \xeCJK_pre_inter_class_toks:nnn {AnalysisLatinLetter} {#1}
          { \group_end: }
      }
  }
% ASCII 字母及常用拉丁扩展字母；明确跳过乘号、除号，不重分类标点。
\xeCJKDeclareCharClass { AnalysisLatinLetter }
  { "0041->"005A, "0061->"007A, "00C0->"00D6, "00D8->"00F6,
    "00F8->"024F, "1E00->"1EFF }
\ExplSyntaxOff
% 数学组内局部关闭字母切换，连公式中 \text{...} 的字体也保持原样。
% 只追加，不覆盖 amsmath、unicode-math、xeCJK 已有的数学初始化。
\AtBeginDocument{%
  \everymath\expandafter{\the\everymath\AnalysisPreserveMathText}%
  \everydisplay\expandafter{\the\everydisplay\AnalysisPreserveMathText}%
}


% 定理类环境（定理、定义、命题、引理、推论、断言、猜想、开放问题、注记、
% 习题）的正文：中文用楷体，西文保持正体，数学字体不受影响。
% 注意 \normalfont 只重置西文字族，中文字族仍随默认（宋体），故须显式 \kaishu；
% 不能退回 \itshape——那会把西文一并变斜。
\newtheoremstyle{analysisplain}
  {1.2ex}{1.2ex}{\normalfont\upshape\kaishu}{}
  {\color{AnalysisBlue}\bfseries}{.}{0.6em}{}
\newtheoremstyle{analysisdefinition}
  {1.2ex}{1.2ex}{\normalfont\kaishu}{}
  {\color{AnalysisBlue}\bfseries}{.}{0.6em}{}
\newtheoremstyle{analysisremark}
  {1.0ex}{1.0ex}{\normalfont\kaishu}{}
  {\color{AnalysisGray}\bfseries}{.}{0.6em}{}

\theoremstyle{analysisplain}
\newtheorem{theorem}{定理}[section]
% 正文默认按“章.节.流水号”编号；未分节的附录省略中间的 0。
\renewcommand{\thetheorem}{%
  \ifnum\value{section}=0
    \thechapter.\arabic{theorem}%
  \else
    \thesection.\arabic{theorem}%
  \fi
}
\newaliascnt{lemma}{theorem}
\newtheorem{lemma}[lemma]{引理}
\aliascntresetthe{lemma}
\newaliascnt{proposition}{theorem}
\newtheorem{proposition}[proposition]{命题}
\aliascntresetthe{proposition}
\newaliascnt{corollary}{theorem}
\newtheorem{corollary}[corollary]{推论}
\aliascntresetthe{corollary}
\newaliascnt{claim}{theorem}
\newtheorem{claim}[claim]{断言}
\aliascntresetthe{claim}
\newaliascnt{conjecture}{theorem}
\newtheorem{conjecture}[conjecture]{猜想}
\aliascntresetthe{conjecture}
\newaliascnt{openproblem}{theorem}
\newtheorem{openproblem}[openproblem]{开放问题}
\aliascntresetthe{openproblem}

\theoremstyle{analysisdefinition}
\newaliascnt{definition}{theorem}
\newtheorem{definition}[definition]{定义}
\aliascntresetthe{definition}
% 习题按章单独编号，不占用定理类环境的流水号。
\newtheorem{exercise}{习题}[chapter]

\theoremstyle{analysisremark}
\newaliascnt{remark}{theorem}
\newtheorem{remark}[remark]{注记}
\aliascntresetthe{remark}
\renewcommand{\proofname}{证明}
% 保留 amsthm 的证明间距与终止符。标题（证明／解／证明概要）的中文用楷体，
% 西文保持正体；正文不施加楷体，仍为默认宋体。不能退回 \itshape。
\makeatletter
\renewenvironment{proof}[1][\proofname]{\par
  \pushQED{\qed}%
  \normalfont \topsep6\p@\@plus6\p@\relax
  \trivlist
  \item[\hskip\labelsep\upshape\kaishu #1\@addpunct{.}]\ignorespaces
}{%
  \popQED\endtrivlist\@endpefalse
}
\makeatother
\newenvironment{solution}
  {\begin{proof}[解]}
  {\end{proof}}
\newenvironment{proofsketch}
  {\begin{proof}[证明概要]}
  {\end{proof}}

% 正文选读内容的星号属于编号而不属于标题；第二参数统一接收标签。
\newcommand{\AnalysisOptionalMark}{\texorpdfstring{\textsuperscript{*}}{*}}
% 选做习题沿用同一编号星号。可选参数仍只表示习题自身的简短题名，
% 不得再写“选读”或“选做”；星号由本环境自动附在题号右侧。
\NewDocumentEnvironment{optionalexercise}{o}{%
  \begingroup
  \let\AnalysisSavedTheExercise\theexercise
  \renewcommand{\theexercise}{\AnalysisSavedTheExercise\AnalysisOptionalMark}%
  \IfNoValueTF{#1}{\begin{exercise}}{\begin{exercise}[#1]}%
}{%
  \end{exercise}%
  \endgroup
}
% 可选参数只供目录与 PDF 书签使用；正文题名可借此按语义手动换行，
% 又不把换行命令写入目录或书签。
\NewDocumentCommand{\OptionalChapter}{O{}mm}{%
  \begingroup
  \let\AnalysisSavedTheChapter\thechapter
  \renewcommand{\thechapter}{\AnalysisSavedTheChapter\AnalysisOptionalMark}%
  \ctexset{chapter={name={第,章\AnalysisOptionalMark},
    number=\chinese{chapter},aftername=\quad}}%
  \ifstrempty{#1}{\chapter{#2}}{\chapter[#1]{#2}}\label{#3}%
  \endgroup
}
\NewDocumentCommand{\OptionalSection}{mm}{%
  \begingroup
  \let\AnalysisSavedTheSection\thesection
  \renewcommand{\thesection}{\AnalysisSavedTheSection\AnalysisOptionalMark}%
  \ctexset{section={number=\thesection}}%
  \section{#1}\label{#2}%
  \endgroup
}
\NewDocumentCommand{\OptionalSubsection}{mm}{%
  \begingroup
  \let\AnalysisSavedTheSubsection\thesubsection
  \renewcommand{\thesubsection}{\AnalysisSavedTheSubsection\AnalysisOptionalMark}%
  \ctexset{subsection={number=\thesubsection}}%
  \subsection{#1}\label{#2}%
  \endgroup
}

% 编题页必须从奇数页开始，题页仅在页脚显示页码，其后保留一张空白页。
\makeatletter
\NewDocumentCommand{\BookPart}{mm}{%
  \cleardoublepage
  \begingroup
  \def\AnalysisPartLabel{#2}%
  % \part 会在返回前翻页，故须在 \@endpart 开始处记录题页标签。
  \pretocmd{\@endpart}{\label{\AnalysisPartLabel}}{}%
    {\PackageError{analysis}{无法为编题页写入标签}{请检查当前文档类的 \string\@endpart 定义。}}%
  \part{#1}%
  \endgroup
  \thispagestyle{empty}%
  \null
  \clearpage
}
% 各编的“本编导读”是星号章（不计编号、不入目录编号）。ctex 的 \@schapter
% 调用 \CTEX@makeanchor@schapter，该宏在 ctex 中只有调用而无定义，展开为空，
% 故星号章不产生 PDF 书签。此处显式补一个章级（第 1 层）书签，使“本编导读”
% 与各章并列挂在所属编之下。
\NewDocumentCommand{\BookPartGuide}{O{本编导读}}{%
  \chapter*{#1}%
  \pdfbookmark[1]{#1}{AnalysisPartGuide@\thepage}%
}

% “附录”只作为书签与目录中的编级节点存在，不排可见的编题页：
% 第一个附录章直接自奇数页开始，PDF 书签中仍保留“附录”顶层节点。
% startatroot 使根级书签与书后顶层节点分离；phantomsection 须先于
% addcontentsline，目录中的“附录”行才能获得独立锚点而非沿用前一节的
% section*.N。根级“附录”书签连同其子章由 addcontentsline 的 part 行自动
% 建立，不得再手写 \pdfbookmark（实测会产生重复的“附录”节点）。
\NewDocumentCommand{\BookAppendixPart}{m}{%
  \cleardoublepage
  \bookmarksetup{startatroot}%
  \phantomsection
  \addcontentsline{toc}{part}{附录}%
  \label{#1}%
}
\makeatother

% 列表项的括号与编号整体承担链接，避免只让中间的编号可点击。
\NewDocumentCommand{\itemref}{m}{\hyperref[#1]{(\ref*{#1})}}

% 章、节、小节的引用：整个“第 X 章／第 X 节／第 X 小节”都是跳转链接。
% 链接必须走 hyperref 的标签解析机制（目标名 chapter.N／section.N.M）；
% 早先用 \hyperlink{裸标签名} 指向不存在的目标，链接全部失效。
%
% 章号用汉字（“第五章”）且必须取“被引用那一章”的编号；\thechapter 是当前
% 计数器，在编导读、前言或跨章引用时会给出错误章号（表现为“第零章”）。
% hyperref 下 \ref 受保护、\getrefnumber 在 \edef 中又失真，故改由 zref 记录：
% 各章在 \chapter 之后调用 \BookChapterNumber，把汉字号存为 zref 属性；
% \zref@extract 可展开，能在 \edef 中安全取出。
\usepackage[user,hyperref]{zref}
\makeatletter
\zref@newprop{analysischapternum}{}
% 在 \chapter 之后调用，参数即该章标签；把本章汉字号存为 zref 属性
\NewDocumentCommand{\BookChapterNumber}{m}{%
  \zref@setcurrent{analysischapternum}{\chinese{chapter}}%
  \zref@labelbyprops{zch@#1}{analysischapternum}%
}
% 取出某章标签对应的汉字号（\zref@extract 可展开，能在 \edef 中安全使用）
\NewDocumentCommand{\AnalysisChineseChapterOf}{m}{%
  \zref@extract{zch@#1}{analysischapternum}}
\makeatother
% \ref* 的分组会遮住编号末端的西文字符类别；编号前的间距可由 xeCJK
% 自动产生，编号后的同类间距须在完整链接内部显式补回。
\NewDocumentCommand{\secref}{m}{%
  \hyperref[#1]{第\ref*{#1}\CJKecglue 节}}
\NewDocumentCommand{\subsecref}{m}{%
  \hyperref[#1]{第\ref*{#1}\CJKecglue 小节}}
\NewDocumentCommand{\chapref}{m}{%
  \begingroup
  \xeCJKsetup{CJKecglue={}}%
  \edef\AnalysisChapterDisp{第\AnalysisChineseChapterOf{#1}章}%
  \hyperref[#1]{\AnalysisChapterDisp}%
  \endgroup}

% “弱-*拓扑”等固定术语在源码中虽已连写，xeCJK 仍会在西文星号与
% 后续汉字之间自动加入中西文胶。该零宽连接记号只抑制这一处自动
% 间距；写入书签时展开为空，避免把排版命令带入 PDF 字符串。
\NewDocumentCommand{\WeakStarJoin}{}{%
  \texorpdfstring{\kern0pt{}}{}}
\pdfstringdefDisableCommands{\def\WeakStarJoin{}}

% 算法采用带编号的步骤，并明确标出输入与输出。
\floatname{algorithm}{算法}
\renewcommand{\listalgorithmname}{算法目录}
\algrenewcommand\algorithmicrequire{\textbf{输入：}}
\algrenewcommand\algorithmicensure{\textbf{输出：}}
\numberwithin{algorithm}{chapter}
\NewDocumentEnvironment{algorithmsteps}{mm}
  {\begin{algorithmic}[1]\Require #1\Ensure #2}
  {\end{algorithmic}}

\newtcolorbox{learninggoals}{
  enhanced,breakable,colback=AnalysisBlueLight,colframe=AnalysisBlue,
  boxrule=.45pt,arc=1mm,left=1.2mm,right=1.2mm,top=1mm,bottom=1mm,
  title=本章目标,fonttitle=\bfseries
}
\newtcolorbox{partguide}{
  enhanced,breakable,colback=AnalysisGrayLight,colframe=AnalysisGray,
  boxrule=.35pt,arc=1mm,left=1.2mm,right=1.2mm,top=1mm,bottom=1mm,
  title=本编导读,fonttitle=\bfseries
}
\newtcolorbox{chaptersummary}{
  enhanced,breakable,colback=AnalysisBlueLight,colframe=AnalysisBlue,
  boxrule=.45pt,arc=1mm,left=1.2mm,right=1.2mm,top=1mm,bottom=1mm,
  title=本章小结,fonttitle=\bfseries,
  % 与阅读提示框一致：仅本框内容为多段散文，须恢复首行缩进。
  before upper={\parindent=2em\parskip=0pt}
}
% 仅用于尚待高阶数学复核的核心断言；成稿前必须处理或移除。
\newtcolorbox{coreverification}{
  enhanced,breakable,colback=yellow!10,colframe=orange!85!black,
  boxrule=.55pt,arc=1mm,left=1.2mm,right=1.2mm,top=1mm,bottom=1mm,
  title=待 Sol Ultra 复核,fonttitle=\bfseries
}
\NewDocumentEnvironment{chapterabstract}{}
  {\list{}{\rightmargin\leftmargin}\item\relax\small\color{AnalysisGray}}
  {\endlist}
% 引诗与短引文统一居中，中文使用定理类环境的楷体；章摘要已在上方独立定义，
% 不继承这里的字形与对齐方式。
\RenewDocumentEnvironment{quote}{}
  {\list{}{\rightmargin\leftmargin}\item\relax\centering\kaishu}
  {\endlist}
% 引诗环境保留为语义清楚的兼容入口，版式与 quote 完全一致。
\NewDocumentEnvironment{poemquote}{}
  {\begin{quote}}
  {\end{quote}}
% 引诗末尾的作者或出处在引诗区域内右对齐。
\NewDocumentCommand{\quoteattribution}{m}{%
  \par\smallskip
  {\raggedleft #1\par}%
}

\makeindex[name=chinese,title=中文索引,columns=2,intoc,options=-q]
\makeindex[name=foreign,title=外文索引,columns=2,intoc,options=-q]
\makeindex[name=symbols,title=符号索引,columns=2,intoc,options=-q]

% 有条目时读入 makeindex 生成的真实索引；尚无条目时仍保留书后结构与目录入口。
% 一键脚本会在 .idx 为空时删除旧 .ind，故这里不会误排上一次构建的陈旧索引。
\NewDocumentCommand{\PrintBookIndex}{m m}{%
  \IfFileExists{#1.ind}{%
    \printindex[#1]%
  }{%
    \chapter*{#2}%
    \addcontentsline{toc}{chapter}{#2}%
  }%
}

% 所有术语在每卷首次正式定义或引入时都显示完整原文，并登记中外文索引。
% 星号形式与缺失原文均视为错误，防止旧的隐藏原文规则被重新带回源码。
% 定理类正文处于楷体；Fandol 楷体没有粗体字形，直接 \textbf 会退回普通楷体。
% 术语主称因此局部恢复正文的宋体字族，再使用与普通正文一致的粗体。
\NewDocumentCommand{\term}{s O{} m o}{%
  \IfBooleanT{#1}{%
    \PackageError{analysis}{\string\term* 已取消}{请改用 \string\term，并保留末尾的原文注解。}%
  }%
  {\songti\textbf{\textup{#3}}}%
  \IfNoValueTF{#4}{%
    \PackageError{analysis}{术语缺少原文注解}{每个术语首次正式定义或引入时都必须写成 \string\term[排序键]{中文主称}[原文]。}%
  }{%
    \textnormal{(\textup{#4})}%
  }%
  \ifstrempty{#2}{\index[chinese]{#3}}{\index[chinese]{#2@#3}}%
  \IfNoValueF{#4}{\index[foreign]{#4}}%
}
\NewDocumentCommand{\termidx}{O{}m}{%
  \ifstrempty{#1}{\index[chinese]{#2}}{\index[chinese]{#1@#2}}%
}
% 中文异名不另排成带页码的同义主词条，而以“见 主称”导向规范名称。
\renewcommand*{\seename}{见}
\NewDocumentCommand{\termsee}{O{}mm}{%
  \ifstrempty{#1}{\index[chinese]{#2|see{#3}}}{\index[chinese]{#1@#2|see{#3}}}%
}
\NewDocumentCommand{\foreignidx}{O{}m}{%
  \ifstrempty{#1}{\index[foreign]{#2}}{\index[foreign]{#1@#2}}%
}
\newcounter{analysisSymbolIndexOrder}
\NewDocumentCommand{\symidx}{m}{%
  \stepcounter{analysisSymbolIndexOrder}%
  \edef\AnalysisSymbolSortKey{%
    \number\numexpr10000+\value{analysisSymbolIndexOrder}\relax}%
  \index[symbols]{\AnalysisSymbolSortKey@#1}%
}

% 一般强调与术语粗体分离。
\NewDocumentCommand{\keypoint}{m}{\CJKunderdot{#1}}
\NewDocumentCommand{\foreignkeypoint}{m}{\underline{\textup{#1}}}

\crefname{chapter}{章}{章}
\crefname{part}{编}{编}
\crefname{section}{节}{节}
\crefname{subsection}{小节}{小节}
\crefname{theorem}{定理}{定理}
\crefname{lemma}{引理}{引理}
\crefname{proposition}{命题}{命题}
\crefname{corollary}{推论}{推论}
\crefname{claim}{断言}{断言}
\crefname{conjecture}{猜想}{猜想}
\crefname{openproblem}{开放问题}{开放问题}
\crefname{definition}{定义}{定义}
\crefname{remark}{注记}{注记}
\crefname{exercise}{习题}{习题}
\crefname{equation}{式}{式}
\crefname{figure}{图}{图}
\crefname{table}{表}{表}
\crefname{algorithm}{算法}{算法}

% cleveref 会在自己的分组中排出“定义3.1.5”一类完整链接；分组结束后，
% xeCJK 无法再从末端识别出西文数字，因而不会在紧随其后的汉字前插入
% 与“义”和“3”之间相同的 \CJKecglue。这里只在下一个源码记号确为
% CJK 正文字时补回同一种间距；空格、控制序列以及中英文标点均不处理，
% 从而不在“定义3.1.5，”或“定义3.1.5.”前制造额外空白。
\NewCommandCopy{\AnalysisOriginalCref}{\cref}
\NewCommandCopy{\AnalysisOriginalCrefCapital}{\Cref}
\ExplSyntaxOn
\cs_new_protected:Npn \analysis_cref_maybe_cjk_glue:
  { \peek_after:Nw \analysis_cref_maybe_cjk_glue_aux: }
\cs_new_protected:Npn \analysis_cref_maybe_cjk_glue_aux:
  {
    \token_if_space:NF \l_peek_token
      {
        \token_if_cs:NF \l_peek_token
          {
            \xeCJK_if_same_class:NNT \l_peek_token 中
              { \CJKecglue }
          }
      }
  }
\RenewDocumentCommand{\cref}{s m}
  {
    \IfBooleanTF{#1}
      { \AnalysisOriginalCref*{#2} }
      { \AnalysisOriginalCref{#2} }
    \analysis_cref_maybe_cjk_glue:
  }
\RenewDocumentCommand{\Cref}{s m}
  {
    \IfBooleanTF{#1}
      { \AnalysisOriginalCrefCapital*{#2} }
      { \AnalysisOriginalCrefCapital{#2} }
    \analysis_cref_maybe_cjk_glue:
  }
% 附录标签沿用字母编号；名称与编号整体进入同一链接，并与 \cref
% 采用相同的链接后中西文间距规则。
\NewDocumentCommand{\appref}{m}
  {
    \hyperref[#1]{附录\CJKecglue\ref*{#1}}
    \analysis_cref_maybe_cjk_glue:
  }
\ExplSyntaxOff

\newcommand{\MakeBookTitle}[1][封面]{%
  % 书名页在 PDF 侧栏中作为顶级入口；卷级文件可传入更具辨识性的书签名。
  \bookmarksetup{startatroot}%
  \pdfbookmark[0]{#1}{AnalysisCover}%
  \begin{titlepage}
    \centering
    \vspace*{2.2cm}
    {\fontsize{36}{46}\selectfont\sffamily 分析学\par}
    \vspace{0.7cm}
    {\Large\sffamily\color{AnalysisGray}\BookVolumeName\par}
    \vspace{1.0cm}
    {\fontsize{24}{32}\selectfont\bfseries\color{AnalysisBlue}\BookTitle\par}
    \vfill
    {\large \BookAuthor\ 编写\par}
    \vspace{0.5cm}
    {\normalsize\BookDraftDate\par}
    \vspace*{1.8cm}
  \end{titlepage}
}

% 目录也应有独立的顶级 PDF 书签。使用公共命令避免三卷分别遗漏。
\newcommand{\MakeBookContents}{%
  \cleardoublepage
  \bookmarksetup{startatroot}%
  \pdfbookmark[0]{目录}{AnalysisContents}%
  \tableofcontents
}

% 某卷首次出现表格后，在总目录之后生成表目录；无表格的卷不调用本命令。
\newcommand{\MakeBookListOfTables}{%
  \cleardoublepage
  \bookmarksetup{startatroot}%
  \pdfbookmark[0]{表目录}{AnalysisListOfTables}%
  \begingroup
    \renewcommand{\listtablename}{表目录}%
    \listoftables
  \endgroup
}

% 无编号章（序言、目录、编导读、后记、参考文献、索引等）同样写入页眉标记，
% 使后续偶数页页眉显示该单元标题。置于 hyperref 之后，保证保留 hyperref 的章锚点。
\makeatletter
\let\Analysis@orig@schapter\@schapter
\def\@schapter#1{%
  \markboth{#1}{}%
  \Analysis@orig@schapter{#1}%
}
% 目录按层级着色：编为黑色、小节为深灰；章与节保持默认链接蓝。
\pretocmd{\l@part}{\hypersetup{linkcolor=black}}{}{%
  \PackageWarning{analysis}{无法为目录编条目着色}}
\pretocmd{\l@chapter}{\hypersetup{linkcolor=AnalysisBlue}}{}{%
  \PackageWarning{analysis}{无法为目录章条目着色}}
\pretocmd{\l@section}{\hypersetup{linkcolor=AnalysisBlue}}{}{%
  \PackageWarning{analysis}{无法为目录节条目着色}}
\pretocmd{\l@subsection}{\hypersetup{linkcolor=AnalysisGray}}{}{%
  \PackageWarning{analysis}{无法为目录小节条目着色}}
% 编的目录条目与 PDF 书签须用纯文本“第X编”。ctex 的 \CTEX@part@tocline
% 取 \CTEXthepart，而该宏带上编题页的颜色与字体设置，会把格式代码写进目录，
% 并使书签出现 “=AnalysisGray第一编” 之类乱码。此覆盖须置于 \ctexset 之后，
% 否则会被 ctex 重新定义；\CTEXifname 是 ctex 的内部开关，此处沿用。
\def\CTEX@part@tocline#1#2{%
  \CTEXifname
    {第\chinese{part}编\hspace{1em}}%
    {}%
  #2}
\makeatother
